\newpage
\section*{Solucionario completo}

\noindent Cada respuesta incluye la clave y una justificación breve. En los ejercicios numéricos se indica la ecuación o el cálculo principal.

\subsection*{Bloque A. Fundamentos teóricos y conceptuales}

\sol{1}{A}{El monopolio es una estructura de mercado con un único oferente, sin sustitutos cercanos y con barreras que impiden la entrada de competidores.}

\sol{2}{B}{El monopolista enfrenta la demanda de mercado, que tiene pendiente negativa. Por ello puede elegir el precio o la cantidad, pero no ambos independientemente.}

\sol{3}{B}{Las barreras a la entrada pueden provenir del control de un insumo esencial, patentes, licencias, economías de escala o ventajas tecnológicas.}

\sol{4}{C}{La empresa competitiva es precio-aceptante y enfrenta una demanda horizontal; el monopolio enfrenta la demanda de mercado y tiene poder para fijar el precio.}

\sol{5}{D}{El monopolio se enfrenta a la curva de demanda de mercado, que es descendente: para vender más debe reducir el precio.}

\sol{6}{B}{Al bajar el precio para vender una unidad adicional, el monopolista pierde ingreso en todas las unidades anteriores; por eso $IMg<P$.}

\sol{7}{C}{La condición de maximización es $IMg=CMg$. El precio se lee luego sobre la curva de demanda.}

\sol{8}{C}{El precio de monopolio se ubica sobre la demanda y por encima del costo marginal, lo que genera el margen monopólico.}

\sol{9}{A}{Las barreras a la entrada impiden que nuevos competidores erosionen los beneficios extraordinarios, por lo que pueden persistir en el largo plazo.}

\sol{10}{C}{El monopolio no tiene una curva de oferta independiente: la cantidad y el precio dependen de la forma de la demanda y no existe una relación única entre precio y costo marginal.}

\sol{11}{D}{El índice de Lerner mide el poder de mercado como $(P-CMg)/P$. Cuanto mayor es, mayor es el margen sobre el costo marginal.}

\sol{12}{D}{La regla de markup es $P=CMg/(1-1/|\varepsilon|)$ para $|\varepsilon|>1$. Relaciona el precio con el costo marginal y la elasticidad.}

\sol{13}{A}{El monopolista produce en el tramo elástico de la demanda, donde $|\varepsilon|>1$; en el tramo inelástico, un aumento de precio aumentaría el ingreso y reduciría el costo.}

\sol{14}{A}{La discriminación de precios consiste en cobrar precios diferentes por el mismo bien a distintos consumidores o unidades, sin que la diferencia se explique solo por costos.}

\sol{15}{D}{En la discriminación de primer grado, el vendedor cobra a cada consumidor su disposición máxima a pagar, apropiándose de todo el excedente del consumidor.}

\sol{16}{B}{La discriminación de segundo grado utiliza autoselección: los consumidores eligen entre opciones como descuentos por cantidad, bloques de consumo o tarifas binomias.}

\sol{17}{C}{La discriminación de tercer grado separa a los consumidores en grupos identificables y cobra precios distintos según sus elasticidades.}

\sol{18}{A}{Se requiere poder de mercado, capacidad de identificar grupos o consumidores y ausencia de arbitraje que permita revender el bien.}

\sol{19}{B}{La discriminación puede aumentar la cantidad y reducir la pérdida de eficiencia, pero también puede extraer excedente y afectar la equidad; el efecto neto es ambiguo.}

\sol{20}{C}{La pérdida de eficiencia del monopolio es el área entre la demanda y el costo marginal para las unidades que no se producen por el precio monopólico.}

\sol{21}{B}{El monopolio fija $P>CMg$ y produce menos que la cantidad eficiente; esa brecha genera la pérdida de bienestar.}

\sol{22}{D}{El monopolio natural presenta costos medios decrecientes en el rango relevante de demanda, lo que hace eficiente que una sola empresa abastezca el mercado.}

\sol{23}{B}{En un monopolio natural, $CMg<CMe$ en el rango relevante. Si se fija $P=CMg$, el precio no cubre el costo medio y la empresa incurre en pérdidas.}

\sol{24}{A}{La regulación puede usar precios máximos, tarifas en dos partes, subsidios o mecanismos de competencia por el mercado para corregir el poder monopólico.}

\sol{25}{C}{La tarifa en dos partes combina un cargo fijo, que puede cubrir costos fijos, con un precio por unidad cercano al costo marginal.}

\sol{26}{D}{El monopolio fija $P>CMg$, por lo que no se cumple la condición de eficiencia asignativa del Primer Teorema del Bienestar.}

\sol{27}{C}{Si está protegido por barreras, el monopolio puede no tener incentivos suficientes para minimizar costos; por eso puede ser productivamente ineficiente.}

\sol{28}{A}{La ineficiencia X se refiere a costos excesivos o gestión ineficiente que surgen por la falta de presión competitiva.}

\sol{29}{D}{Las políticas incluyen la aplicación de la política de competencia, la regulación tarifaria, las subastas de concesión y la propiedad pública.}

\sol{30}{C}{Un precio máximo igual al costo medio permite recuperar costos y evitar el precio monopólico, aunque no alcanza la eficiencia de primer mejor asociada a $P=CMg$.}

\subsection*{Bloque B. Aplicaciones cuantitativas y optimización}

\sol{31}{C}{$IT=(100-Q)Q=100Q-Q^2$, por lo que $IMg=100-2Q$. Igualando $IMg=CMg=20$: $100-2Q=20$, de donde $Q=40$.}

\sol{32}{B}{Sustituyendo $Q=40$ en la demanda: $P=100-40=60$.}

\sol{33}{C}{El beneficio es $(P-CMg)Q=(60-20)40=1600$. También $IT-CT=2400-800=1600$.}

\sol{34}{B}{En el óptimo, $IMg=CMg=20$.}

\sol{35}{A}{$L=(P-CMg)/P=(60-20)/60=2/3\approx 0.67$.}

\sol{36}{D}{Con $Q=100-P$, la elasticidad es $\varepsilon=-(P/Q)=-60/40=-1.5$. En valor absoluto, 1.5.}

\sol{37}{B}{$EC=0.5(100-60)(40)=800$.}

\sol{38}{D}{La cantidad competitiva es 80. $PCP=0.5(80-40)(60-20)=800$.}

\sol{39}{B}{Con $P=CMg=20$: $100-Q=20$, de donde $Q=80$.}

\sol{40}{C}{El precio competitivo es $P=CMg=20$.}

\sol{41}{D}{El excedente total competitivo es $0.5(100-20)(80)=3200$.}

\sol{42}{A}{En el monopolio, $ET=EC+\pi=800+1600=2400$.}

\sol{43}{B}{Con discriminación de primer grado, el monopolista captura todo el excedente del consumidor; por tanto, $EC=0$.}

\sol{44}{A}{El monopolista captura todo el excedente total, que con $Q=80$ es $0.5(100-20)(80)=3200$.}

\sol{45}{B}{Con discriminación perfecta, produce hasta que $P=CMg$ en la última unidad: $Q=80$.}

\sol{46}{A}{En el mercado 1, $IMg_1=100-2Q_1=20$, de donde $Q_1=40$.}

\sol{47}{B}{En el mercado 2, $IMg_2=80-2Q_2=20$, de donde $Q_2=30$.}

\sol{48}{D}{El precio en el mercado 2 es $P_2=80-30=50$.}

\sol{49}{A}{El grupo con demanda más elástica es más sensible al precio; el monopolista le cobra un precio menor.}

\sol{50}{B}{El grupo con demanda menos elástica tiene menor sensibilidad al precio; el monopolista le cobra un precio mayor.}

\sol{51}{B}{$P=CMg/(1-1/|\varepsilon|)=10/(1-1/2)=20$.}

\sol{52}{D}{$P=30/(1-1/4)=30/(3/4)=40$.}

\sol{53}{C}{$L=(50-30)/50=0.4$.}

\sol{54}{A}{$L=(P-CMg)/P=0.5$. Con $P=80$: $80-CMg=40$, de donde $CMg=40$.}

\sol{55}{C}{$IMg=50-2Q$. Igualando $IMg=CMg=10$: $50-2Q=10$, de donde $Q=20$.}

\sol{56}{D}{$P=50-20=30$.}

\sol{57}{D}{$CMe=100/20+10=15$. Beneficio $=(30-15)20=300$.}

\sol{58}{B}{Con $P=CMg=10$: $Q=40$. $CMe=100/40+10=12.5$. Pérdida $=(10-12.5)40=-100$.}

\sol{59}{D}{Con precio igual al costo medio, el ingreso por unidad iguala el costo medio y el beneficio económico es cero.}

\sol{60}{C}{En el monopolio natural no regulado, el precio de monopolio es 30 y el costo medio en $Q=20$ es 15; por tanto, $P>CMe$.}

\sol{61}{A}{Como $P=CMg<CMe$, la empresa pierde; se requiere un subsidio u otro mecanismo para cubrir la diferencia.}

\sol{62}{D}{El cargo fijo permite cubrir los costos fijos y el precio variable puede acercarse al costo marginal, reduciendo la distorsión.}

\sol{63}{C}{El precio de monopolio es 60 y el competitivo es 20; la diferencia es 40.}

\sol{64}{A}{El beneficio del monopolista es $(60-20)40=1600$.}

\sol{65}{B}{Cuanto menos elástica es la demanda, mayor es el margen precio-costo y mayor tiende a ser la pérdida de eficiencia.}

\subsection*{Bloque C. Casos con implicancias regulatorias y tarifarias}

\sol{66}{A}{Una empresa de agua con costos medios decrecientes en el rango relevante es un monopolio natural; una sola empresa puede abastecer el mercado a menor costo que varias.}

\sol{67}{C}{En un monopolio natural, $P=CMg<CMe$ genera pérdidas; se necesita un subsidio, una tarifa en dos partes u otro mecanismo de cobertura.}

\sol{68}{D}{El monopolio no regulado fija $P>CMg$ y produce menos que la cantidad eficiente, lo que genera la pérdida de bienestar.}

\sol{69}{C}{Con $P=CMe$, la empresa recupera costos y obtiene beneficio económico cero; es una solución de segundo mejor.}

\sol{70}{B}{La tarifa en dos partes combina un cargo fijo y un precio variable por unidad consumida.}

\sol{71}{B}{La discriminación de tercer grado separa grupos identificables, como usuarios residenciales, comerciales e industriales.}

\sol{72}{D}{Se requiere poder de mercado, capacidad de separar grupos y ausencia de arbitraje para que los grupos no revendan el servicio.}

\sol{73}{C}{El grupo con demanda más inelástica soporta un precio mayor; el grupo más elástico recibe un precio menor.}

\sol{74}{D}{La discriminación de primer grado captura todo el excedente del consumidor; puede aumentar la producción, pero plantea problemas distributivos.}

\sol{75}{A}{El monopolio natural se define por costos medios decrecientes en el rango relevante de demanda, lo que lo convierte en un caso de subaditividad de costos.}

\sol{76}{A}{La regulación por tasa de retorno puede incentivar la sobreinversión, porque la empresa obtiene retorno sobre el capital invertido.}

\sol{77}{C}{Un price cap fija un límite al precio y deja que la empresa se apropie de las reducciones de costos, creando incentivos a la eficiencia.}

\sol{78}{D}{La subasta de concesión introduce competencia por el mercado y puede disciplinar al operador aunque exista un solo prestador.}

\sol{79}{D}{El yardstick competition compara el desempeño de empresas similares y usa esa referencia para fijar estándares o tarifas.}

\sol{80}{B}{La producción pública o municipal puede alinear objetivos sociales, pero también enfrenta riesgos de presupuesto blando y menor presión competitiva.}

\sol{81}{B}{La regulación, la competencia y las subastas pueden reducir el margen monopólico y la pérdida de eficiencia.}

\sol{82}{A}{El subsidio a hogares vulnerables se justifica por equidad, aunque puede distorsionar el consumo y requiere financiamiento.}

\sol{83}{B}{La tarifa en bloques crecientes es una forma de discriminación de segundo grado: el precio depende del volumen consumido y el usuario se autoselecciona.}

\sol{84}{D}{La tarifa en bloques crecientes puede promover la conservación y proteger consumos básicos, pero debe equilibrarse con la eficiencia y la sostenibilidad financiera.}

\sol{85}{A}{El cargo fijo permite cubrir costos fijos y reducir la distorsión que generaría un precio variable muy alto.}

\sol{86}{C}{La regulación enfrenta un trade-off entre eficiencia asignativa y sostenibilidad financiera: $P=CMg$ es eficiente pero puede generar pérdidas.}

\sol{87}{A}{Sin regulación, el monopolio fija $P>CMg$ y produce menos que la cantidad eficiente.}

\sol{88}{A}{La solución de primer mejor es $P=CMg$; en un monopolio natural esto requiere un subsidio que cubra la pérdida.}

\sol{89}{C}{La solución de segundo mejor suele ser $P=CMe$ o una estructura de precios Ramsey, que equilibra eficiencia y cobertura de costos.}

\sol{90}{C}{La discriminación de precios puede mejorar la recuperación de costos, pero también puede afectar la equidad y requiere evaluación regulatoria.}

\subsection*{Bloque D. Ampliación: Teoría y Aplicaciones Prácticas (91 a 180)}

\sol{91}{A}{Dado $IMg = P(1 + 1/\varepsilon)$, si $|\varepsilon| = 0.8 < 1$, entonces $IMg < 0 < CMg$. Reduciendo la cantidad siempre aumenta el ingreso total y disminuye el costo total, por lo que no existe óptimo interior finito.}
\sol{92}{B}{De la condición de primer orden $IMg = CMg$ y la relación de Amoroso-Robinson $IMg = P(1 + 1/\varepsilon)$, se deduce $L = (P - CMg)/P = -1/\varepsilon = 1/|\varepsilon|$.}
\sol{93}{C}{Con demanda lineal $P = a - bQ$, el precio óptimo es $P^* = (a + c + t)/2$, por lo que $\partial P^*/\partial t = 1/2$. Exactamente la mitad del impuesto se traslada al consumidor.}
\sol{94}{D}{Con elasticidad constante $|\varepsilon| > 1$, el precio es $P^* = \frac{c+t}{1 - 1/|\varepsilon|}$. Entonces $\frac{\partial P^*}{\partial t} = \frac{1}{1 - 1/|\varepsilon|} > 1$, produciéndose sobretraslado.}
\sol{95}{A}{Al maximizar $\pi = P(Q_1+Q_2)(Q_1+Q_2) - C_1(Q_1) - C_2(Q_2)$, las CPO exigen $IMg(Q_1+Q_2) = CMg_1(Q_1) = CMg_2(Q_2)$.}
\sol{96}{B}{Sin compromiso, tras vender a los consumidores de alta valoración, el monopolista tiene incentivos a bajar el precio para atender a la demanda residual. Anticipando esto, los compradores esperan y el poder de mercado colapsa hacia $P = CMg$.}
\sol{97}{C}{El alquiler (leasing) elimina la competencia con unidades ya vendidas en el mercado secundario, y la cláusula de nación más favorecida hace costoso bajar precios futuros porque obliga a reembolsar la diferencia a compradores previos.}
\sol{98}{D}{Harberger mide solo el triángulo de pérdida de eficiencia asignativa, mientras que Tullock señala que la competencia por capturar el rectángulo de beneficios monopólicos (lobby, litigios) puede disipar hasta la totalidad de esas rentas.}
\sol{99}{A}{El monopolista valora la calidad según el consumidor marginal (que fija el precio), mientras el planificador social la valora según el consumidor promedio. Si el marginal valora menos la calidad que el promedio, el monopolio subprovee calidad.}
\sol{100}{B}{La condición de segundo orden exige $\pi''(Q^*) = IMg'(Q^*) - CMg'(Q^*) < 0$, es decir, $IMg'(Q^*) < CMg'(Q^*)$: el costo marginal debe cortar al ingreso marginal desde abajo.}
\sol{101}{C}{Al cobrar a cada unidad su disposición máxima a pagar, $IMg(Q) = P(Q)$, por lo que el monopolista produce hasta $P(Q) = CMg(Q)$ (eficiencia asignativa, $PBE = 0$), apropiándose de la totalidad del excedente.}
\sol{102}{D}{En el óptimo $IMg_A = IMg_B = CMg$, lo que implica $P_A(1 - 1/|\varepsilon_A|) = P_B(1 - 1/|\varepsilon_B|)$. Como $|\varepsilon_A| < |\varepsilon_B|$, el término entre paréntesis es menor en A, por lo que $P_A > P_B$.}
\sol{103}{A}{Varian (1985) demostró que $\Delta W \le (P_u - c)\Delta Q$ (cota superior con precios diferenciados). Si $\Delta Q \le 0$, el bienestar social necesariamente cae por ineficiencia en la asignación entre consumidores.}
\sol{104}{B}{Con demandas lineales y ambos mercados atendidos bajo precio uniforme, la suma de cantidades vendidas no cambia ($\Delta Q = 0$). Al existir precios distintos entre mercados, se pierde eficiencia en el consumo y $\Delta W < 0$.}
\sol{105}{C}{En el óptimo con autoselección (Mussa-Rosen / Maskin-Riley), el tipo alto consume la cantidad eficiente ($P = CMg$) y retiene renta informacional para que no imite al tipo bajo; el tipo bajo sufre distorsión hacia abajo y obtiene excedente cero.}
\sol{106}{D}{En el óptimo, el monopolista extrae todo el excedente del tipo bajo ($IR_L$ activa) y ofrece al tipo alto la renta justa para que no elija el paquete del tipo bajo ($IC_H$ activa).}
\sol{107}{A}{Para discriminar en tercer grado se requiere: (1) poder de mercado, (2) capacidad de identificar grupos con elasticidades diferentes y (3) imposibilidad o alto costo de arbitraje/reventa entre grupos.}
\sol{108}{B}{Las tarifas por bloques crecientes cobran precios menores (incluso por debajo del costo medio) en el bloque social/básico y precios más altos en bloques superiores, promoviendo uso racional y financiamiento cruzado.}
\sol{109}{C}{Degradar la versión baja reduce su atractivo para los consumidores de alta valoración, relajando la restricción $IC_H$ y permitiendo al monopolista cobrar un precio mayor por la versión superior.}
\sol{110}{D}{Si los consumidores de alta disposición a pagar tienen alta paciencia (factor de descuento cercano a 1) y anticipan la caída del precio, pospondrán su compra, frustrando la separación intertemporal.}
\sol{111}{A}{En tercer grado el vendedor usa características observables (edad, categoría comercial/doméstica) para fijar precios lineales distintos; en segundo grado no observa el tipo y diseña menús para inducir autoselección.}
\sol{112}{B}{Cuando el precio marginal cae en bloques sucesivos, el gasto total es cóncavo en $Q$, por lo que el precio medio $G(Q)/Q$ decrece con $Q$ y se sitúa por encima del precio marginal del último tramo.}
\sol{113}{C}{Si bajo precio único el mercado débil no es atendido ($Q_2 = 0$), discriminar permite abrir ese mercado sin alterar el precio en el mercado fuerte (con $CMg$ constante), generando excedente puro adicional.}
\sol{114}{D}{Los cupones funcionan como un mecanismo de autoselección (screening): quienes tienen menor costo de oportunidad del tiempo y demanda más elástica invierten esfuerzo en usarlos y pagan un precio menor.}
\sol{115}{A}{El precio uniforme maximizador de beneficios sobre la demanda agregada satisface $P_B < P_u < P_A$ (donde A es el mercado menos elástico y B el más elástico), siempre que ambos sigan activos.}
\sol{116}{B}{Con consumidores homogéneos, fijar $p = c$ maximiza el excedente total de cada cliente, y fijar $T = EC(c)$ permite al monopolista capturar el 100\% de ese excedente con eficiencia asignativa.}
\sol{117}{C}{Al subir $p$ ligeramente sobre $CMg$, la pérdida de cargo fijo $T$ (fijado por el tipo bajo) es compensada en primer orden por el mayor margen $(p - CMg)$ cobrado sobre el volumen adicional del tipo alto.}
\sol{118}{D}{La correlación negativa reduce la dispersión (varianza) de la valoración del paquete frente a las valoraciones individuales, permitiendo extraer una mayor fracción del excedente total.}
\sol{119}{A}{Con $CMg > 0$, el empaquetamiento puro obliga a entregar ambos bienes incluso a quien valora uno por debajo de su $CMg$. El mixto permite que ese cliente compre solo el bien que valora por encima del costo.}
\sol{120}{B}{El bien complementario (cartuchos) mide la intensidad de uso del bien principal (impresora), permitiendo cobrar más a los usuarios intensivos de alta disposición a pagar sin observar su tipo.}
\sol{121}{C}{Con proporciones fijas y demanda competitiva de B, el precio que el consumidor paga por el sistema es $P_A + P_B$. El monopolista ya extrae el máximo beneficio fijando $P_A$ con $P_B = CMg_B$ (teorema de la renta única).}
\sol{122}{D}{En el caso de pico firme, solo los usuarios del periodo pico presionan la expansión de capacidad, por lo que $P_{\text{pico}} = b + \beta$, mientras los usuarios en periodo valle solo generan costo operativo $P_{\text{valle}} = b$.}
\sol{123}{A}{En el caso de pico cambiante, ambos periodos utilizan la capacidad completa ($Q_1 = Q_2 = K$) y contribuyen a financiarla: $P_1 > b$, $P_2 > b$ con $(P_1 - b) + (P_2 - b) = \beta$, es decir, $P_1 + P_2 = 2b + \beta$.}
\sol{124}{B}{Cobrar el verdadero costo marginal (incluyendo capacidad) en horas pico desincentiva el consumo en congestión y traslada uso hacia horas valle, optimizando la capacidad instalada.}
\sol{125}{C}{Con consumidores homogéneos, la tarifa en dos partes fija $p = CMg$, eliminando el triángulo de Harberger ($PBE = 0$), y extrae todo el excedente mediante el cargo fijo $T$.}
\sol{126}{D}{Si atiende a ambos, $T$ está acotado por el excedente del grupo bajo; si excluye al grupo bajo, fija $p = CMg$ y $T = EC_H(CMg)$, lo que es más rentable si hay suficientes clientes de tipo alto con excedente muy superior.}
\sol{127}{A}{Un cargo fijo demasiado alto actúa como barrera de acceso (restricción de participación) para hogares de bajos ingresos, lo que motiva que reguladores como SUNASS mantengan cargos fijos bajos.}
\sol{128}{B}{Whinston demostró que si hay economías de escala en el bien B y el atado es un compromiso irreversible, el monopolista reduce las ventas viables del rival por debajo de la escala mínima eficiente, disuadiendo su entrada.}
\sol{129}{C}{Baumol (1977) estableció que la subaditividad de costos ($C(Q) < \sum C(q_i)$ para toda partición no trivial $\sum q_i = Q$) es la condición exacta de monopolio natural.}
\sol{130}{D}{Las economías de escala son suficientes pero no necesarias para la subaditividad. Si $Q$ supera ligeramente la escala mínima eficiente $Q_{min}$, dividir $Q$ entre dos firmas puede seguir siendo más caro que producirlo en una sola.}
\sol{131}{A}{En producción multiproducto, ni las economías de escala radiales ni las economías de alcance por sí solas garantizan subaditividad; se requiere su combinación (o convexidad trans-radial y economías de escala).}
\sol{132}{B}{Si la tarifa regulada excede el costo en segmentos de bajo costo o alta rentabilidad (para financiar subsidios cruzados o costos fijos), un entrante puede captar solo esos clientes rentables y desfinanciar al monopolio natural.}
\sol{133}{C}{En un mercado perfectamente contestable (sin costos hundidos ni rezagos de respuesta), la amenaza de entrada «golpea y huye» (hit-and-run) obliga al incumbente a fijar precios sostenibles con beneficio económico cero.}
\sol{134}{D}{Cuando $CMe$ es decreciente, $CMg < CMe$. Por tanto, el precio de primer mejor ($P = CMg$) implica ingresos menores a los costos totales ($PQ < CT$), generando déficit operativo que debe cubrirse con subsidios o tarifas de segundo mejor.}
\sol{135}{A}{La regla de Ramsey-Boiteux es $\frac{P_i - CMg_i}{P_i} = \frac{\lambda}{1+\lambda}\frac{1}{|\varepsilon_i|}$: para minimizar la pérdida de eficiencia al cubrir costos fijos, se carga un mayor margen donde la demanda es menos elástica.}
\sol{136}{B}{Los servicios básicos esenciales tienen demandas muy inelásticas; cargar mayores márgenes sobre ellos resulta regresivo. Por ello, los reguladores suelen introducir ponderaciones distributivas (Ramsey con equidad) o subsidios focalizados.}
\sol{137}{C}{Cuando $\lambda \to \infty$ (restricción presupuestaria muy severa), $R = \lambda/(1+\lambda) \to 1$, y la regla se convierte en $(P_i - CMg_i)/P_i = 1/|\varepsilon_i|$, que es exactamente la regla de Lerner del monopolio.}
\sol{138}{D}{Una asignación de ingresos está libre de subsidios cruzados si $C(q_1, q_2) - C(q_j) \le R_i \le C(q_i)$: ningún servicio paga menos que su costo incremental ni más de lo que costaría producirlo por separado.}
\sol{139}{A}{Si $R_i > C(q_i)$ (el costo stand-alone), entonces como $R_1 + R_2 = C(q_1, q_2)$, el otro servicio paga $R_j < C(q_1, q_2) - C(q_i)$ (menos que su costo incremental), por lo que $i$ subsidia a $j$.}
\sol{140}{B}{Con un tope $P_{\max}$, el monopolista no pierde ingresos en unidades previas al expandir ventas hasta $Q(P_{\max})$; así su $IMg = P_{\max} > CMg$ y aumenta la producción hasta $Q(P_{\max})$.}
\sol{141}{C}{Existen economías de alcance si $C(q_1, q_2) < C(q_1, 0) + C(0, q_2)$, es decir, producir conjuntamente ahorra costos frente a la producción especializada en dos empresas distintas.}
\sol{142}{D}{La regla ECPR fija $a = c_{\text{acceso}} + (P_{\text{final}} - c_{\text{minorista}} - c_{\text{acceso}}) = P_{\text{final}} - c_{\text{minorista}}$, garantizando que solo entren rivales al menos tan eficientes como el incumbente aguas abajo.}
\sol{143}{A}{Como el beneficio máximo permitido es $(s - r)K$, la firma puede ampliar sus utilidades expandiendo la base de capital $K$, lo que reduce el costo sombra del capital y genera sobrecapitalización ($TMST_{L,K}$ distorsionada).}
\sol{144}{B}{Al desvincular temporalmente la tarifa de los costos realizados durante el quinquenio regulatorio, la empresa se apropia de las reducciones de costos que logre por encima del factor $X$, incentivando eficiencia productiva.}
\sol{145}{C}{Bajo price cap puro, la firma puede recortar calidad para bajar costos (requiriendo metas de gestión) y, si anticipa que mayores ahorros hoy elevarán el factor $X$ futuro (efecto trinquete), moderará su esfuerzo.}
\sol{146}{D}{Al vincular el precio de la firma $i$ al costo promedio de las demás firmas similares ($j \neq i$), la firma $i$ no puede influir en su propia tarifa inflando costos, replicando incentivos competitivos.}
\sol{147}{A}{Si la firma recibe $PQ + EC(P) - C(Q) = W(Q)$ (el bienestar social total), al maximizar sus utilidades elegirá exactamente $P = CMg$. Su limitación es el alto costo fiscal y las grandes rentas transferidas a la firma.}
\sol{148}{B}{Dar contratos de alto poder (precio fijo) maximiza el esfuerzo pero deja grandes rentas informacionales a firmas eficientes; contratos de bajo poder (costo-plus) reducen rentas pero destruyen incentivos al esfuerzo.}
\sol{149}{C}{Análogo a la discriminación de segundo grado («no distorsión en el tope»), la firma de bajo costo (eficiente) elige el contrato de precio fijo (esfuerzo óptimo, renta positiva) y la ineficiente elige un contrato con reparto de costos.}
\sol{150}{D}{Demsetz propuso licitar la concesión por la menor tarifa (no por el mayor canon, que solo transferiría rentas monopólicas al Estado), logrando en teoría $P = CMe$ sin regulación continua.}
\sol{151}{A}{SUNASS utiliza un modelo de empresa regulada con horizonte quinquenal basado en el PMO (metas de gestión e inversiones), asegurando que el valor presente de ingresos cubra costos eficientes y remuneración del capital al WACC.}
\sol{152}{B}{El reajuste por IPM protege el equilibrio económico-financiero en términos reales ante variaciones del nivel general de precios sin necesidad de abrir un nuevo estudio tarifario.}
\sol{153}{C}{Los subsidios cruzados generalizados por volumen subsidian incluso a hogares no pobres de bajo consumo; la focalización reduce filtraciones (error de inclusión) y mejora la cobertura a hogares pobres (error de exclusión).}
\sol{154}{D}{El ingreso marginal es $IMg = 120 - 4Q$ y el costo marginal es $CMg = 20$. Igualando: $120 - 4Q = 20 \Rightarrow Q^* = 25$, $P^* = 120 - 50 = 70$ y $\pi^* = 70(25) - (20(25) + 100) = 1750 - 600 = 1150$.}
\sol{155}{A}{En monopolio $100 - 4Q = 20 \Rightarrow Q_m = 20, P_m = 60$. En competencia $100 - 2Q = 20 \Rightarrow Q_c = 40, P_c = 20$. La pérdida de eficiencia es $PBE = \frac{1}{2}(60 - 20)(40 - 20) = 400$.}
\sol{156}{B}{De $L = \frac{P - CMg}{P} = 0.40$, se obtiene $0.60 P = 30 \Rightarrow P^* = 50.0$. Además, $L = -1/\varepsilon \Rightarrow \varepsilon = -1/0.40 = -2.50$.}
\sol{157}{C}{Aplicando la regla $P^* = \frac{CMg}{1 + 1/\varepsilon} = \frac{18}{1 - 1/2.5} = \frac{18}{0.6} = 30.00$.}
\sol{158}{D}{Sin impuesto, $100 - 2Q = 20 \Rightarrow Q_0 = 40, P_0 = 60$. Con $t = 12$, el nuevo costo marginal es $32$, por lo que $100 - 2Q = 32 \Rightarrow Q_1 = 34, P_1^* = 66$, resultando $\Delta P = 6 = t/2$.}
\sol{159}{A}{Con $\varepsilon = -3$, $P^* = \frac{CMg}{1 - 1/3} = 1.5 \cdot CMg$. Así, $P_0 = 30$ y con $CMg_1 = 26$, $P_1 = 39$, por lo que $\Delta P = 1.5(6) = 9.00$ (sobretraslado del 150\%).}
\sol{160}{B}{Igualando costos marginales: $2Q_1 = 4Q_2 \Rightarrow Q_1 = 2Q_2$, luego $Q = 3Q_2$ y $CMg(Q) = \frac{4}{3}Q$. Igualando a $IMg = 100 - 2Q$: $100 = \frac{10}{3}Q \Rightarrow Q = 30$, con $Q_1^* = 20$, $Q_2^* = 10$ y $P^* = 70$.}
\sol{161}{C}{Tenemos $IMg = 100 - 2Q$ y $CMg = Q + 10$. Igualando: $3Q = 90 \Rightarrow Q^* = 30$ y $P^* = 70$. El beneficio es $\pi^* = 70(30) - [0.5(900) + 10(30) + 100] = 2100 - 850 = 1250$.}
\sol{162}{D}{En primer grado produce hasta $P = CMg$: $80 - 2Q = 20 \Rightarrow Q^* = 30$. El ingreso total equivale al área bajo la demanda: $\frac{80 + 20}{2}(30) = 1500$. Restando $CT = 20(30) + 100 = 700$, resulta $\pi^* = 800$.}
\sol{163}{A}{Igualando $IMg_i = CMg$ en cada mercado: $100 - 2Q_1 = 20 \Rightarrow Q_1^* = 40, P_1^* = 60$; y $60 - 2Q_2 = 20 \Rightarrow Q_2^* = 20, P_2^* = 40$.}
\sol{164}{B}{De la regla de Lerner en cada segmento: $P_1^* = \frac{15}{1 - 1/2} = 30$ y $P_2^* = \frac{15}{1 - 1/4} = 20$. Por tanto, $P_1^* / P_2^* = 30 / 20 = 1.50$.}
\sol{165}{C}{Bajo precio único, $Q = 160 - 2P \Rightarrow P = 80 - 0.5Q$, $IMg = 80 - Q = 20 \Rightarrow Q = 60, P = 50$ y $\pi_u = (50 - 20)(60) = 1800$. Discriminando, $\pi_d = (60 - 20)(40) + (40 - 20)(20) = 2000$. La ganancia neta es $200$.}
\sol{166}{D}{Según el modelo de Oi, se fija $p^* = CMg = 10$, con lo cual $q^* = (50 - 10)/0.5 = 80$. El cargo fijo extrae todo el excedente del consumidor: $T^* = \frac{1}{2}(50 - 10)(80) = 1600$.}
\sol{167}{A}{Fijando $T = EC_L(p) = \frac{1}{2}(12-p)^2$, el beneficio es $\pi(p) = 3\cdot\frac{1}{2}(12-p)^2 + (p-4)(42-3p)$. Derivando: $\pi'(p) = -3(12-p) + (42-3p) - 3(p-4) = 18 - 3p = 0 \Rightarrow p^* = 6$, $T^* = 18$ y $\pi = 102 > 98$ (excluir).}
\sol{168}{B}{Por separado, fijar $P_A = P_B = 100$ genera $100 + 100 = 200$ (mejor que $40 \times 2 = 80$ por bien). Con paquete puro, ambos valoran el paquete en $140$, logrando $140 \times 2 = 280$.}
\sol{169}{C}{El consumidor 1 compra solo A ($90 - 20 = 70$), el consumidor 2 compra el paquete ($120 - 40 = 80$) y el consumidor 3 compra solo B ($90 - 20 = 70$). El beneficio total es $70 + 80 + 70 = 220$, superior al paquete puro ($180$).}
\sol{170}{D}{En el caso de pico firme, $P_1^* = b + \beta = 30 \Rightarrow Q_1^* = 70$ y $P_2^* = b = 10 \Rightarrow Q_2^* = 30$. Como $Q_2^* < Q_1^*$, la solución es consistente y la capacidad óptima es $K^* = 70$.}
\sol{171}{A}{Con $P = CMg = 10$, se obtiene $Q^* = 90 - 10 = 80$. El ingreso total es $10(80) = 800$ y el costo total es $400 + 10(80) = 1200$, generando un déficit exactamente igual al costo fijo de $400$.}
\sol{172}{B}{Igualando $60 - Q = \frac{400}{Q} + 10$ se obtiene $Q^2 - 50Q + 400 = 0$, cuyas raíces son $Q = 10$ y $Q = 40$. Para maximizar bienestar se elige la raíz mayor: $Q^* = 40$ y $P^* = 20$.}
\sol{173}{C}{Para dos firmas simétricas, $C(Q) < 2C(Q/2) \iff 50 + 2Q^2 < 100 + Q^2 \iff Q < \sqrt{50} \approx 7.07$. En $Q=6$, $C(6)=122 < 2C(3)=136$ (subaditiva); en $Q=8$, $C(8)=178 > 2C(4)=164$ (no subaditiva).}
\sol{174}{D}{De $\frac{P_i - CMg_i}{P_i} = \frac{R}{|\varepsilon_i|}$: para el bien 1, $\frac{P_1 - 12}{P_1} = \frac{0.3}{1.5} = 0.2 \Rightarrow P_1^* = 15.00$; para el bien 2, $\frac{P_2 - 9}{P_2} = \frac{0.3}{3.0} = 0.1 \Rightarrow P_2^* = 10.00$.}
\sol{175}{A}{Aplicando la fórmula lineal de ajuste por precio tope: $P_t = 4.00 \times (1 + 0.05 - 0.02) = 4.00 \times 1.03 = S/\ 4.12/\text{m}^3$.}
\sol{176}{B}{El ingreso requerido es $IR = OPEX + \text{Depreciación} + sK = 400 + 100 + 0.08(2000) = 660$. Dividiendo entre el volumen $Q = 220$, la tarifa media es $660 / 220 = 3.00$.}
\sol{177}{C}{Según Faulhaber, $R_1$ debe estar entre su costo incremental $C(q_1, q_2) - C(0, q_2) = 150 - 80 = 70$ y su costo stand-alone $C(q_1, 0) = 100$, es decir, $70 \le R_1 \le 100$.}
\sol{178}{D}{La factura se compone del cargo fijo ($S/\ 5.00$), los primeros $10\ \text{m}^3$ a $S/\ 1.50$ ($S/\ 15.00$) y los siguientes $10\ \text{m}^3$ a $S/\ 3.00$ ($S/\ 30.00$), totalizando $S/\ 50.00$.}
\sol{179}{A}{Sin regulación, $Q_m = 40, P_m = 60$ y $PBE_m = \frac{1}{2}(40)(40) = 800$. Con $P_{\max} = 40$, el monopolista produce $Q_{\text{reg}} = 60$ y $PBE_{\text{reg}} = \frac{1}{2}(20)(20) = 200$, reduciendo la PBE en $600$.}
\sol{180}{B}{La condición de Dorfman-Steiner establece que $\frac{A}{PQ} = \frac{\varepsilon_A}{|\varepsilon_P|}$. Sustituyendo: $\frac{0.5}{2.5} = 0.20$, es decir, el $20\%$ de los ingresos totales.}
